% ==============================================================================
% 01_einleitung.tex – Kapitel 1
% ==============================================================================
\chapter{Einleitung}
\label{ch:einleitung}

\section{Motivation}
\label{sec:motivation}
Eye-Tracking misst, wohin ein Mensch schaut. Diese Information ist für viele
Anwendungen nützlich: für die Bedienung von Computern ohne Hände, zum
Beispiel für Menschen mit körperlichen Einschränkungen, für Studien zur
Benutzerfreundlichkeit von Webseiten, für Lernplattformen oder für die
Forschung in Psychologie und Medizin. Professionelle Eye-Tracker arbeiten mit
Infrarotlicht und speziellen Kameras. Sie sind sehr genau, aber teuer und
nicht überall verfügbar.

Eine normale Webcam ist dagegen in fast jedem Laptop eingebaut. Wenn man den
Blick mit einer Webcam schätzen könnte, wäre Eye-Tracking für sehr viele
Menschen ohne zusätzliche Hardware möglich. Das ist aber schwierig. Die Bilder
haben eine niedrige Auflösung, und die Augen sind darin nur wenige Pixel groß.
Das Licht ändert sich, der Kopf bewegt sich, viele Menschen tragen eine Brille,
und jede Person hat eine andere Augen- und Gesichtsform. Ein allgemeines Modell,
das auf vielen Personen trainiert wurde, passt deshalb oft schlecht zu einer
neuen Person.

\section{Problemstellung und Ziel}
\label{sec:ziel}
Ein Ausweg ist die \emph{Personalisierung}: Man nimmt ein vortrainiertes
Modell und passt es mit wenigen Daten an eine Person an. Die Person schaut dazu
kurz auf einen Punkt, der sich über den Bildschirm bewegt (Kalibrierung). Die
Kalibrierung soll aber kurz sein, denn niemand möchte vor jeder Nutzung
mehrere Minuten lang kalibrieren. Die Personalisierung muss also
\emph{dateneffizient} sein.

Das Ziel dieser Arbeit ist es, zu untersuchen, wie gut ein vortrainiertes
Bildverarbeitungsmodell mit einer kurzen Webcam-Kalibrierung an eine Person
angepasst werden kann. Als Modell dient ein auf ImageNet vortrainiertes
ResNet-18. Sein Backbone bleibt eingefroren, und nur ein kleiner
Regression Head wird trainiert. Der Head wird zuerst auf dem Datensatz
GazeCapture trainiert und danach mit eigenen Aufnahmen personalisiert.

\section{Forschungsfragen}
\label{sec:forschungsfragen}
Die Arbeit beantwortet vier Forschungsfragen (FF):
\begin{description}[leftmargin=1.2cm, style=nextline]
  \item[FF1] Wie stark senkt eine kurze Personalisierung den Fehler gegenüber
        einem allgemeinen Modell und gegenüber einer einfachen Schätzung, die
        immer denselben Punkt ausgibt?
  \item[FF2] Welche Trainingseinstellung eignet sich für den Regression Head
        (trainierbare Schichten, Activation, Loss Function, Learning Rate,
        Batch Size)?
  \item[FF3] Wie beeinflussen die Aufteilung der Daten (zufällig, zeitlich,
        zwei Videos), die Gestaltung der Kalibrierung (Bewegung des Ziels,
        Entfernen von Transition Frames) und eine zeitliche Glättung der
        Vorhersagen das Ergebnis?
  \item[FF4] Welchen Einfluss haben die Aufnahmebedingungen (Licht,
        Hintergrund, Brille, Dauer, Gerät bzw.\ Webcam)?
\end{description}

\section{Beitrag der Arbeit}
\label{sec:beitrag}
Die Arbeit liefert folgende Beiträge:
\begin{itemize}
  \item \textbf{Eigene Aufnahmen:} 27 Kalibriersitzungen von Person~A (dem
        Autor) mit zwei Laptop-Webcams, davon 18 ausgewertet, und zwei
        Sitzungen einer zweiten Person (Person~B) als Vergleich. Die Sitzungen unterscheiden sich in der
        Bewegung des Ziels, im Licht, im Hintergrund, in der Brille und in der
        Dauer. Die neueste Generation speichert für jedes
        Bild die Phase des Ziels (\code{WAIT} oder \code{MOVE}).
  \item \textbf{Reproduzierbare Pipeline:} feste Seeds, eingefrorener Backbone,
        gespeicherte Features (Feature-Cache) und gespeicherte Vorhersagen.
        Ein Personalisierungslauf dauert damit etwa zwei Sekunden auf einer CPU.
  \item \textbf{Korrigiertes Fehlermaß:} Ein Fehler im Auswertungscode wurde
        gefunden und behoben. Baselines werden immer auf denselben Daten wie
        das Modell berechnet.
  \item \textbf{Systematische Studie:} Hyperparameter mit fünf Wiederholungen
        und Confidence Intervals, räumliche Analyse des Fehlers, ein
        Vergleich von Datenaufteilungen und Kalibriergestaltungen sowie eine
        Untersuchung der zeitlichen Glättung der Vorhersagen.
  \item \textbf{Web-Demo:} Eine Demo im Browser zeigt die Personalisierung
        live: Kalibrierung, Blickverfolgung und Auswertung (Kapitel~6).
\end{itemize}

\section{Aufbau der Arbeit}
\label{sec:aufbau}
Kapitel~\ref{ch:sota} beschreibt den Stand der Technik.
Kapitel~\ref{ch:methodik} erklärt Daten, Modell und Vorgehen.
Kapitel~\ref{ch:results} zeigt die Experimente und Ergebnisse.
Kapitel~\ref{ch:diskussion} beantwortet die Forschungsfragen und nennt die
Grenzen der Arbeit. Kapitel~\ref{ch:fazit} fasst die Arbeit zusammen und gibt
einen Ausblick.

In dieser Arbeit werden die üblichen englischen Fachbegriffe verwendet, zum
Beispiel Overfitting, Feature oder Loss. Ihre deutsche Übersetzung steht im
Verzeichnis der Fachbegriffe im Anhang (\cref{sec:anhang_fachbegriffe}).
